\section{Conditioning on avoidance of one orientation}

Let \(B_n\) count connected meanders with no positive site, and put
\(b=\mu-1\), so \(B_n^{1/n}\to b\). Give this avoidance class the
uniform probability measure. Write \(Y\) for the number of negative
sites under this measure. In this section the site set is
\(V_n=\{0,\ldots,n-2\}\).

\begin{theorem}\label{thm:conditional-ratio}
The one-orientation avoidance counts satisfy
\[
             \lim_{n\to\infty}\frac{B_n}{B_{n-1}}=b .
\]
In the interior, the negative-site indicators converge to a stationary
two-dependent process. For a finite set \(T\subset\mathbb Z\), with
\(k=|T|\) and \(d=|((T-1)\cup(T+1))\setminus T|\), its intersection
probabilities are
\[
   \mathbb P(\text{all sites of }T\text{ are negative})
             =b^{-k}(1+b^{-1})^d .
\]
In particular its one-site density equals \((b+1)^2/b^3\).
\end{theorem}

\begin{proof}
At finite order fix \(T\subseteq V_n\) and take the exterior neighbours
only in \(V_n\), so that
\(d=|(((T-1)\cup(T+1))\cap V_n)\setminus T|\).
A negative site excludes a positive site at its own location and at
each neighbour. There remain \(L=n-1-k-d\) sites at which a positive
occurrence must be excluded. Fixed-site contraction and
inclusion--exclusion give the count
\[
 \sum_{j=0}^{L}(-1)^j\binom Lj M_{n-k-j}
       =\sum_{a=0}^{d}\binom da B_{n-k-a}.                 \tag{*}
\]
For the equality, the left-hand sum also counts order-\((n-k)\)
meanders avoiding positives at any prescribed \(L\) sites.
Their exact positive set is a subset of the other \(d\) sites.
The exact-set formula gives the right-hand sum.

The ratio follows from Theorem~\ref{thm:idle-shifts} at \(t=1\).
Divide \((*)\) by \(B_n\) and use the consecutive-ratio limit for
each of its finitely many terms. Sets \(T\) translated into the
interior retain both exterior neighbours whenever these are distinct.
The resulting intersection probabilities determine every finite law by
inclusion--exclusion and are translation invariant. If two sets have
mutual distance at least three, both \(k\) and \(d\) add for each choice
of subsets; their joint law factors. Singletons at distance two have
three distinct exterior neighbours in total, giving strictly less than
the product of their one-site probabilities.
\end{proof}

The same law arises from conditioning in a finite window of the
unconditioned limiting coloured process. Let \(p=1/\mu\), let \(T\) be
finite, and let \(W\) contain \(T\) and its neighbours. Prescribing
negatives on \(T\) excludes positives on \(k+d\) sites of \(W\).
Inclusion--exclusion on the remaining sites gives the intersection
probability \(p^k(1-p)^{|W|-k-d}\). Dividing by
\(\mathbb P(\text{no positive in }W)=(1-p)^{|W|}\) gives exactly
\(b^{-k}(1+b^{-1})^d\), independently of how far \(W\) extends beyond
the neighbours. Thus conditioning the finite meander on avoidance
throughout its full site set has the same local limit as this
finite-window calculation.

\subsection*{A finite entropy comparison in the full class}
For \(n\ge2\) put \(N=n-1\), \(p=1/\mu\), and \(q_n=M_{n-1}/M_n\).
Writing \(D\) for binary relative entropy, one has
\[
 N D(q_n\Vert p)\le
      \log\!\left(\frac{b\mu^{\,N}}{M_n}\right),\qquad
 |q_n-p|\le
      \sqrt{\frac{1}{2N}\log\!\left(\frac{b\mu^{\,N}}{M_n}\right)}.
\]
Indeed \(B_j\le b^j\) by supermultiplicativity. A particular binary
indicator word of weight \(k\) has probability \(B_{n-k}/M_n\),
whereas its probability under the product Bernoulli-\(p\) law is
\(b^{N-k}/\mu^N\). Thus the relative entropy of the full law with
respect to the product law is at most the logarithm displayed above.
Subadditivity of Shannon entropy bounds it below by the sum of the
\(N\) one-site divergences. Finally \(D(q\Vert p)\ge2(q-p)^2\),
as follows by twice differentiating in \(q\). The right-hand side of
the first inequality is \(o(n)\) by root growth. This also proves the
unconditioned ratio theorem directly. A numerical error bound from
this expression requires information about \(\mu\).
